\documentclass[11pt]{article}
\usepackage[a4paper,margin=25mm]{geometry}
\usepackage[no-math]{fontspec}
\setmainfont{Latin Modern Roman}
\usepackage{amsmath,amssymb,amsthm,xfrac,xspace,hyperref,graphicx,comment}
\excludecomment{DefTralics}
\providecommand{\C}{}
\newcommand{\ie}{i.e.,\xspace}
\newcommand{\eg}{e.g.,\xspace}
\newcommand{\etc}{etc.\xspace}
\newcommand{\cf}{cf.\xspace}
\newcommand{\resp}{resp.\xspace}
\newcommand{\smaller}{\small}
\newcommand{\Subsection}{\subsection}
\newcommand{\Subsubsection}{\subsubsection}
\newcommand{\skpt}{\smallskip}
\emergencystretch=3em

% Begin literal retained input author-preamble.tex
\usepackage{mathrsfs}
\usepackage{eucal}
\let\mathcal\EuScript
\usepackage[all,cmtip]{xy}\let\labelstyle\textstyle
\newdir^{ (}{{}*!/-5pt/\dir^{(}}
\newdir_{ (}{{}*!/-5pt/\dir_{(}}
\multlinegap0pt

\renewcommand{\theenumi}{\roman{enumi}}
\newcommand\mto{\mathchoice{\longmapsto}{\mapsto}{\mapsto}{\mapsto}}
\newcommand\hto{\mathrel{\lhook\joinrel\to}}
\newcommand\implique{\mathchoice{\implies}{\Rightarrow}{\Rightarrow}{\Rightarrow}}
\newcommand{\To}[1]{\mathchoice{\xrightarrow{\textstyle\kern4pt#1\kern3pt}}{\stackrel{#1}{\longrightarrow}}{}{}}
\newcommand{\isom}{\stackrel{\sim}{\longrightarrow}}
\let\oldbigsqcup\bigsqcup
\renewcommand{\bigsqcup}{\mathop{\textstyle\oldbigsqcup}\displaylimits}
\let\oldvee\vee
\renewcommand{\vee}{{\scriptscriptstyle\oldvee}}
\newcommand{\lpr}{(\!(}
\newcommand{\rpr}{)\!)}
\newcommand{\lcr}{[\![}
\newcommand{\rcr}{]\!]}
\newcommand{\wtj}{\widetilde\jmath}

%\usepackage{xr}
%\externaldocument[BW1-]{partie_1}

\theoremstyle{plain}
\newtheorem{thm}{Theorem}[section]
\newtheorem{thmintro}{Theorem}
\newtheorem{sublem}[thm]{Sublemma}
\newtheorem{lem}[thm]{Lemma}
\newtheorem{conj}[thm]{Conjecture}
\newtheorem{question}[thm]{Question}
\newtheorem{cor}[thm]{Corollary}
\newtheorem{prop}[thm]{Proposition}

\theoremstyle{definition}
\newtheorem{defn}[thm]{Definition}
\newtheorem{rmk}[thm]{Remark}
\newtheorem{rmks}[thm]{Remarks}
\newtheorem{example}[thm]{Example}
\newtheorem{examples}[thm]{Examples}
\newtheorem{Step}{Step}

\numberwithin{equation}{section}

\newcommand{\alg}{{\mathrm{alg}}}
\newcommand{\cl}{{\mathrm{cl}}}
\newcommand{\cx}{{\mathrm{cx}}}
\newcommand{\rs}{{\mathrm{r}}}
\newcommand{\gr}{{\mathrm{gr}}}
\newcommand{\Gm}{\mathbf{G}_\mathrm{m}}
\newcommand{\sB}{{\mathscr B}}
\newcommand{\sH}{{\mathscr H}}
\newcommand{\sF}{{\mathscr F}}
\newcommand{\sG}{{\mathscr G}}
\newcommand{\sI}{{\mathscr I}}
\newcommand{\sO}{{\mathscr O}}
\newcommand{\sX}{{\mathscr X}}
\newcommand{\A}{{\mathbf A}}
\renewcommand{\C}{{\mathbf C}}
\renewcommand{\P}{{\mathbf P}}
\newcommand{\Q}{{\mathbf Q}}
\newcommand{\R}{{\mathbf R}}
\newcommand{\SB}{{\mathbf{SB}}}
\newcommand{\Z}{{\mathbf Z}}
\newcommand{\Zl}{{\Z_{\ell}}}
\newcommand{\CH}{\mathrm{CH}}
\newcommand{\Hdg}{\mathrm{Hdg}}
\newcommand{\Gal}{\mathrm{Gal}}
\newcommand{\nr}{\mathrm{nr}}
\newcommand{\KS}{\textsc{KS}}
\newcommand{\Pic}{\mathrm{Pic}}
\newcommand{\Br}{\mathrm{Br}}
\newcommand{\Tr}{\mathrm{Tr}}
\newcommand{\Id}{\mathrm{Id}}
\newcommand{\Trf}{\mathrm{Tr}_f}
\newcommand{\Spec}{\mathrm{Spec}}
\renewcommand{\phi}{\varphi}
\renewcommand{\emptyset}{\varnothing}
\newcommand{\red}{{\mathrm{red}}}
\newcommand{\sm}{\mathrm{sm}}
\newcommand{\Ker}{{\mathrm{Ker}}}
\newcommand{\Ima}{{\mathrm{Im}}}
\newcommand{\Hom}{{\mathrm{Hom}}}
\newcommand{\Ext}{{\mathrm{Ext}}}
\newcommand{\et}{\text{ét}}
\newcommand{\torsion}{(\mathrm{torsion})}
\newcommand{\tors}{{\mathrm{tors}}}
\newcommand{\elw}{\mathrm{ind}}
\newcommand{\ab}{\mathrm{ab}}
\newcommand{\van}{\mathrm{van}}
\newcommand{\Sing}{\mathrm{Sing}}
\newcommand{\bS}{{\mathbf S}}
\newcommand{\divi}{\mathrm{div}}
\newcommand{\ci}{\mathscr{C}^{\infty}}
\newcommand{\RR}{{\mathrm{R}}}
\hyphenation{semi-stable}
\hyphenation{sub-manifold}
\hyphenation{sub-manifolds}
\hyphenation{mani-fold}
\hyphenation{mani-folds}
\hyphenation{ar-chi-me-dean}


% End literal retained input author-preamble.tex

\begin{document}
\begin{center}{\Large Stable item 1926}\end{center}
\paragraph{Source and attribution.}
Olivier Benoist, Olivier Wittenberg, \emph{On the integral Hodge conjecture for real varieties, II}, Journal de l’École polytechnique — Mathématiques 7 (2020), publisher version, DOI 10.5802/jep.120. \href{https://jep.centre-mersenne.org/articles/10.5802/jep.120/}{Publisher article}. Authors' text: \href{https://creativecommons.org/licenses/by/4.0/}{CC BY 4.0}.
\paragraph{Editorial problem boundary.}
Prove only Proposition 9.17 below, the explicit failure of the EPT property on a K3 surface over the real Puiseux field \(R=\bigcup_n\mathbb R((t^{1/n}))\), with the usual positive-infinitesimal ordering. The selected surface, double-cover equation, strict-transform line bundle, vanishing Borel--Haefliger class, and full obstruction to a real-point-free divisor representative occur in the original proof. Definition 9.15 of EPT is retained. No assertion over arbitrary non-archimedean fields or integral-Hodge-conjecture application is selected; the rest of this article is conservative author context only.
\paragraph{Difficulty (editorial): H3.}
An explicit branched double-cover K3 construction converts a hypothetical real-point-free divisor into an odd-bidegree invariant section. Normalization over a Puiseux valuation ring and specialization of real spectra force an odd-degree real curve on which the defining polynomial must have an impossible constant sign.
\paragraph{Editorial transcription note.}
Original author language, mathematical content and ordering are unchanged. Publisher front matter is replaced by this independent layout wrapper. Original native body and macros remain editable; bibliography is recovered from the matching cached publisher HTML. Adjacent claims are author context, not additional corpus items. Printed mathematical slips are retained without assistant repair.
\clearpage
\paragraph{Standard background (editorial).}
The Borel--Haefliger class of a real line bundle and basic real-spectrum/normalisation facts are standard prerequisites; the class-map setup is discussed in the authors' first article \cite{BW1}, and the exact real-spectrum specialisation fact used in the proof is cited by the authors to \cite[Prop.\,7.1.21]{bcr}. No proof from the first article is asserted or required as a new core step of this selected construction.
\setcounter{section}{9}\setcounter{subsection}{3}\setcounter{subsubsection}{0}\setcounter{thm}{14}

% Begin literal retained input author-definition.tex
\Subsubsection{The signs and EPT properties}
\label{subsubsec:signsEPT}

\begin{defn}
Let $X$ be a smooth projective variety over a real closed field $R$. We say that
\begin{enumerate}
\item $X$ satisfies the \textit{signs property} if for every $K\subset X(R)$ that is a union of connected components, there is a rational function $g$ on~$X$ that is invertible on $X(R)$, such that $g>0$ on $K$ and $g<0$ on $X(R)\setminus K$;
\item $X$ satisfies the \textit{EPT property} if for every line bundle $\mathcal{L}$ on $X$ with vanishing Borel-Haefliger class $\cl_R(\mathcal{L})\in H^1(X(R),\Z/2\Z)$, there exists a divisor $D$ on $X$ such that $\mathcal{L}\simeq\sO_X(D)$ and the support of $D$ contains no $R$-points.
\end{enumerate}
\end{defn}

These two properties are closely related, and complement each other. Together, they completely describe, for any open subset $U\subset X$, the possible distributions of signs $U(R)\to \Z/2\Z$ induced by rational functions $g$ on~$X$ invertible on $U(R)$.

% End literal retained input author-definition.tex

\setcounter{thm}{16}\setcounter{equation}{3}

% Begin literal retained input author-proof.tex
\begin{prop}
\label{prop:cexEPT}
There exists a K3 surface $S$ over $R:=\bigcup_n\R\lpr t^{1/n}\rpr$ that does not satisfy the EPT property.
\end{prop}

\begin{proof}
Consider $\P^1_R\times \P^1_R$ with homogeneous coordinates $([x:y],[v:w])$.
Let $T$ be the double cover of $\P^1_R\times\P^1_R$, with branch locus of bidegree $(4,4)$, defined by the equation
\[
z^2=x(y-x)(xw^2-yv^2)(xw^2+yv^2)-tw^4y^4.
\]
Since $T$ has only rational double points as singularities, its minimal resolution of singularities $S$ is a K3 surface over~$R$. We denote by $i:T\to T$ the involution associated with the double cover.

The unique connected component $K$ of $S(R)$ contained in $0<x/y<1$ is semi-algebraically isomorphic to a sphere, hence satisfies $H^1(K,\Z/2\Z)=0$. The strict transform $\bar{D}$ of $D:=\{v=0\}$ in $S$ has real points only in $K$, so that its Borel-Haefliger class is trivial. Suppose for contradiction that $S$ satisfies the EPT property. Then $\bar{D}$ is linearly equivalent to a divisor whose support has no $R$-points. Pushing forward to $T$ shows that $D$ is linearly equivalent in $T$ to a Weil divisor whose support contains only finitely many $R$-points.

Write $D=E_1-E_2+\divi(f)$, where $E_1$ and $E_2$ are effective Weil divisors whose supports contain finitely many $R$-points. The divisor $E_2+i(E_2)$ comes from $\P^1\times \P^1$, hence is the zero locus of a section in $H^0(T,\sO(d_1,d_2))$ for some $d_1, d_2\in\Z$. It follows that $E_1+i(E_2)$ is the zero locus of a section in $H^0(T,\sO(d_1,d_2+1))$. These two divisors have finitely many $R$-points in their support. Consequently, after possibly replacing~$d_2$ by $d_2+1$, we have constructed $s\in H^0(T,\sO(d_1,d_2))$ with $d_2$ odd such that $\{s=0\}$ contains finitely many $R$-points.

Suppose there exists $a\in T(R)$ such that $s(a)\neq 0$, $i^*s(a)\neq 0$ but $s(a)+i^*s(a)=0$. Choose a semi-algebraic path $\gamma:[0,1]\to T(R)$ joining $a$ and a point $b\in T(R)$ in the ramification locus. Since $s$ vanishes at only finitely many $R$-points, it is possible to assume (after changing $b$) that neither $s$ nor $i^*s$ vanishes along the path. Since $b$ is a ramification point, $s(b)=i^*s(b)$. By hypothesis, $s(a)=-i^*s(a)$.
Since $[0,1]$ is semi-algebraically connected,
the continuous semi-algebraic map
\[
\lambda\mto i^*s(\gamma(\lambda))/s(\gamma(\lambda))
\]
must reach the value $0$, a contradiction. We have proven that $s(a)+i^*s(a)=0$ implies that $s(a)=0$ or $i^*s(a)=0$, hence that $s+i^*s$ vanishes at only finitely many $R$-points of $T$.

Since $s+i^*s$ is $i^*$-invariant, it comes from a
section $s'\in H^0(\P^1_R\times\P^1_R,\sO(d_1,d_2))$ that vanishes at only finitely many
$R$\nobreakdash-points outside of the closed semi-algebraic subset
$\Theta \subset \P^1(R)\times \P^1(R)$
defined by the inequality
\begin{equation}
\label{eq:theta inequality}
x(y-x)(xw^2-yv^2)(xw^2+yv^2)\leq tw^4y^4\rlap{.}
\end{equation}
Replacing $B:=\{s'=0\}$ by an appropriate reduced irreducible component, we
may assume that $B$ is integral.

Let $A=\R\lcr t^{1/n}\rcr$, with~$n$ large enough that
$B=\sB\otimes_A R$ for some closed integral
subscheme $\sB \subset \P^1_A \times \P^1_A$.
Let~$\sB'$ be the normalisation of~$\sB$ and $B'=\sB'\otimes_AR$.
The map $\sB'\otimes_A \R \to \P^1_\R$ induced by the first projection
is generically finite of degree~$d_2$, which is odd;
therefore $\sB' \otimes_A \R$ possesses an irreducible component~$Y$ of odd multiplicity~$m$,
which dominates~$\P^1_\R$ with odd degree.
Letting $Y_\rs \subset \sB'_\rs \supset B'_\rs$ denote the real spectra of~$Y$, $\sB'$
and~$B'$,
the inequality~\eqref{eq:theta inequality} now defines a closed subset~$\Theta_A$ of~$\sB'_\rs$
which contains $B'_\rs$.
The completion of the local ring of~$\sB'$ at the generic point of~$Y$ is isomorphic,
over~$A$, to $\R(Y)\lcr u\rcr$, where $u=(\alpha t^{1/n})^{1/m}$ for some $\alpha \in \R(Y)^*$.
As~$m$ is odd,
any ordering of~$\R(Y)$ can be extended to $\R(Y)\lpr u\rpr$, compatibly with the given ordering on~$A$,
by declaring that $u$ is infinitely small of the same sign as~$\alpha$.
Thus, any point of~$Y_\rs$ above the generic point of~$Y$ belongs to the closure of~$B'_\rs$ in~$\sB'_\rs$
(see \cite[Prop.\,7.1.21]{bcr}) and, hence, to~$\Theta_A$.
Letting~$Y'$ be the normalisation of~$Y$, it follows that the image of~$Y'_\rs$ in~$Y_\rs$
is contained in~$\Theta_A$.

We have now constructed a smooth, proper, irreducible curve~$Y'$ over~$\R$, endowed with two rational
functions $x,v \in \R(Y')^*$, such that the map $x:Y'\to \P^1_\R$ has odd degree and such that
the function
\[
g=x(1-x)(x-v^2)(x+v^2)
\]
takes negative values on~$Y'(\R)$
wherever it is invertible.
There exists $P \in Y'(\R)$ at which~$x$ vanishes to odd order.
The function~$g$ then has a zero or pole of odd order at~$P$; its
sign therefore changes around~$P$, a contradiction.
\end{proof}

% End literal retained input author-proof.tex

\begin{thebibliography}{99}
\bibitem[BCR98]{bcr} J. Bochnak, M. Coste \& M.-F. Roy - Real algebraic geometry , Ergeb. Math. Grenzgeb. (3) , vol. 36 , Springer-Verlag, Berlin, 1998
\bibitem[BW18]{BW1} O. Benoist \& O. Wittenberg - “On the integral Hodge conjecture for real varieties, I” , 2018, à paraître dans Invent. math.
\end{thebibliography}
\end{document}
